\begin{textsample}{POS Dim 4 – 2009\_c5 – Score -1.00 – 2009\_c5/enade\_2009\_biblio\_0e14e9ac55.txt}  \label{ex:f4_pos_260}
Capítulo 5 Distribuição dos Conceitos 5.1 Panorama nacional da distribuição dos conceitos A Tabela 5.1 apresenta a quantidade e o percentual de cursos de BIBLIOTECONOMIA participantes do ENADE/2009 que se situaram em cada uma das faixas de conceitos, distribuídos de acordo com a Região em que se localizam os cursos das instituições de educação superior, mantenedoras dos cursos responsáveis pela oferta.
O Gráfico 5.1 apresenta a mesma informação sob a forma de diagrama de caixa para uma melhor visualização da dispersão dos conceitos obtidos.
No diagrama de caixa, a caixa propriamente dita fica limitada pelos quartis.
A barra dentro da caixa em negrito representa a mediana.
Linhas verticais unem os quartis aos valores extremos.
Se estes valores extremos se configurarem como valores atípicos, estes aparecem como pontos isolados e as linhas terminam no maior ( ou menor ) valor ainda não considerado atípico.
Eventuais valores atípicos aparecem como pequenos círculos com a Categoria Administrativa correspondente.
Valores considerados muito atípicos foram representados por estrelas.
Um diagrama de caixa para poucos dados ou com dados muito concentrados pode ter alguns pontos coincidentes.
Por exemplo, se mais de metade dos cursos de uma Região receberem conceito 5, o valor extremo superior, o quartil superior e a mediana coincidirão.
Se o conjunto só contém um dado este apareceria denotado como a barra da mediana.
Observando -se os dados da Tabela 5.1, percebe -se que, em termos nacionais, predominaram os cursos que receberam os conceitos 3 e 2, correspondentes respectivamente a 39,58 \% e a 31,6 \% deste total.
Ficaram sem conceito um total de 4 cursos ( 10,5 \% ).
Foram classificados com conceitos 1 e 4 outros 7,9 \% cada.
Apenas 1 curso recebeu conceito máximo.
Dos 4 cursos sem conceito, 2 estavam na Região Nordeste e 2 na Região Sudeste.
Nenhum curso apareceu com um valor atípico, muito provavelmente devido ao pequeno número de cursos considerados.
Considerando -se separadamente as cinco Regiões brasileiras, observou -se que, da Região Norte, 2 cursos foram conceituados com conceito 2.
Este aparece no Gráfico 5.1 como uma barra no nível 2.
Na Região Nordeste, 11 cursos participaram da avaliação, 5 com conceito 3 e 4 com conceito 2.
Na Região Sudeste, dos 16 cursos avaliados, 7 cursos, equivalentes a 43,8 \% do total regional, classificaram -se com conceito 3, 4 cursos ( 25,0 \% ) com conceito 2, 2 cursos ( 12,5 \% ) com conceito mínimo e 1 curso ( 6,3 \% ), com conceito 4.
Nenhum curso foi classificado com conceito máximo.
Na Região Sul, foram avaliados 5 cursos dos quais 3 com conceito 3, 1 com conceito 4 e 1 com conceito máximo, 5.
Na Região Centro-Oeste, foram avaliados 4 cursos, 2 com conceito 2, 1 com conceito mínimo e 1 com conceito 4. 5.2 Conceitos por categoria administrativa e por Região A Tabela 5.2, que apresenta a distribuição dos cursos participantes do ENADE/2009 de BIBLIOTECONOMIA, por categoria administrativa, de acordo com os conceitos por eles alcançados, segundo as Regiões percebe -se que, no universo dos 4 cursos que não obtiveram conceito, 3 eram de instituições federais e 1 de instituição privada.
Na Região Nordeste, 2 ficaram sem conceito, ambos de instituição federal.
Na Região Sudeste, 2 cursos não foram avaliados : 1 de instituição federal e 1 de instituição privada.
Nas Regiões Norte, Sul e Centro-Oeste, nenhum ficou sem conceito.
A análise a seguir trata apenas da distribuição dos conceitos dos 34 cursos que foram avaliados, sendo que 20 em instituições federais, 4 em estaduais e 10 em privadas.
O Gráfico 5.2 apresenta a mesma informação sob a forma de diagrama de caixa para uma melhor visualização da dispersão dos conceitos obtidos.
Somente eventuais valores atípicos aparecem com o identificador da Região na qual o curso foi habilitado.
Instituições municipais não estavam representadas nesta área.
De acordo com as informações apresentadas, 58,8 \% dos cursos eram vinculados a instituições federais, 29,4 \% a instituições privadas e 11,8 \% a instituições estaduais.
Apenas 1 curso de instituição federal apresentou conceito máximo, 1 curso, o conceito mínimo e 1 curso, o conceito 4, equivalentes a 5,0 \% do total dessa categoria cada um.
Ainda nas instituições federais, 9 cursos receberam o conceito 3 ( 45,0 \% ) e 8 cursos o conceito 2 ( 40,0 \% ).
Entre as instituições estaduais avaliadas, 2 cursos receberam conceito 3 ( 50,0 \% ), 1 curso recebeu conceito 2 e 1 curso conceito 4 ( 25,0 \% cada ).
Nas instituições privadas 4 cursos receberam conceito 3 ( 40,0 \% ), 3 cursos conceito 2 ( 30,0 \% ) e 1 curso conceito 4 ( 10,0 \% ).
Nenhum curso neste tipo de instituição recebeu conceito máximo e 2 cursos receberam conceito mínimo ( 20,0 \% ).
Considerando -se a análise por Região observou -se que as instituições federais estavam sempre representadas entre as avaliadas.
As instituições estaduais avaliadas estavam presentes apenas nas Regiões Nordeste, Sudeste e Região Sul.
As instituições privadas estavam presentes em todas as Regiões com exceção da Região Norte.
Na Região Norte os 2 cursos avaliados eram de instituição federal e receberam, ambos, o conceito 2.
Na Região Nordeste, dos 9 cursos avaliados, 7 eram de instituição federal, 1 de estadual e 1 de instituição privada.
Dos 7 cursos em instituição federal, 4 receberam conceito 3 e 3 o conceito 2.
O único curso avaliado em instituição estadual recebeu conceito 3 e o único curso avaliado em instituição privada recebeu conceito 2.
Na Região Sudeste dos 14 cursos avaliados, 5 eram de instituição federal, 1 de estadual e 8 de instituições privadas.
Dos 5 cursos em instituição federal, 3 receberam conceito 3, 1 o conceito 2 e 1 o conceito mínimo.
O único curso avaliado em instituição estadual recebeu conceito 2.
Dos 8 cursos avaliados em instituições privadas, 4 receberam conceito 3, 2 o conceito 2 e 1 o conceito 4.
Apenas 1 curso de instituição privada recebeu conceito mínimo e nenhum curso o conceito máximo.
Na Região Sul dos 5 cursos avaliados, 3 eram de instituição federal e 2 de estadual.
Entre os 3 cursos em instituição federal, 2 receberam conceito 3 e 1 o conceito máximo.
Os 2 cursos em instituição estadual receberam os conceitos 3 e 4.
Na Região Centro-Oeste dos 4 cursos avaliados, 3 eram de instituições federais e 1 de instituição privada.
Dos 3 cursos em instituições federais, 2 receberam conceito 2 e 1 o conceito 4.
O único curso avaliado em instituição privada recebeu conceito mínimo. 5.3 Conceitos por Organização Acadêmica e por Região Na Tabela 5.3 encontra -se a distribuição dos conceitos atribuídos aos cursos participantes do ENADE/2009 de BIBLIOTECONOMIA, por Organização Acadêmica, segundo as Regiões.
O Gráfico 5.3 apresenta a mesma informação sob a forma de diagrama de caixa para uma melhor visualização da dispersão dos conceitos obtidos.
Somente eventuais valores atípicos apareceram com o identificador da Categoria Administrativa.
Os cursos desta área estão principalmente alocados em Universidades ( 78,9 \% ) em Faculdades ( 15,8 \% ), depois em Centros Universitários ( 5,3 \% ).
Na Região Norte apenas as Universidades participaram, na Região Nordeste não houve participação de Centros Universitários e na Região Sul, apenas Universidades participaram.
Não houve participação de Centros Federais de Educação Tecnológica ( CEFET ) e Institutos Federais de Educação, Ciência e Tecnologia ( IFECT ) em nenhuma das 5 Regiões.
Nas Universidades e Faculdades, o conceito modal foi 3 ( 43,3 \% e 33,3 \% respectivamente ), nos Centros Universitários o conceito modal foi 2 ( 100 \% ).
Apenas 1 curso em Faculdade e 2 em Universidades obtiveram conceito mínimo e apenas 1 curso em Universidade obteve conceito máximo.
Em 2 tipos de Organização Acadêmica Universidades e Faculdades existiram cursos que ficaram sem conceito ( respectivamente 10,0 \% e 16,7 \% ).
Considerando -se separadamente as cinco Regiões, observa -se que, na Região Norte, apenas 2 cursos em Universidades foram avaliados e obtiveram conceito 2.
Na Região Nordeste, dos 9 cursos avaliados, 8 eram de Universidades, sendo que 5 cursos obtiveram conceito 3 e 3 cursos o conceito 2.
Apenas um curso foi avaliado em Faculdade e obteve conceito 2.
Na Região Sudeste, dos 14 cursos avaliados, 9 era em Universidades, 2 cursos de Centros Universitários e 3 cursos de Faculdades.
Nas Universidades, 5 cursos classificaram -se com conceito 3, 2 cursos com conceito 2 e 2 cursos com conceito mínimo.
Nenhum curso foi classificado nos demais conceitos.
Os 2 cursos de Centros Universitários receberam conceito 2.
Nas Faculdades, 2 cursos receberam conceito 3 e 1 curso o conceito 4.
Na Região Sul, os 5 cursos avaliados, todos em Universidades, classificaram -se com conceito 3 ( 3 cursos ), conceito 4 ( 1 curso ) e conceito máximo ( 1 curso ).
Na Região Centro-Oeste, dos 4 cursos avaliados, 3 eram de Universidades e 1 de Faculdade.
Nas Universidades, 2 cursos receberam conceito 2 e 1 curso, conceito 4.
O único curso em faculdade recebeu conceito mínimo.

% matched lemmas: 
\end{textsample}
